\documentclass{article}
\usepackage[utf8]{inputenc}
\usepackage[T1]{fontenc}
\usepackage{geometry}
\geometry{
    a4paper,
    left=2.5cm,
    right=2.5cm,
    top=2.5cm,
    bottom=2.5cm
}
\usepackage{graphicx}
\usepackage{float}
\usepackage{hyperref}
\usepackage[english]{babel}
\usepackage{amsmath}
\usepackage{amssymb}
\usepackage{array}
\usepackage{booktabs}
\usepackage{caption}
\usepackage{subcaption}
\usepackage{enumitem}
\usepackage{fancyhdr}
\usepackage{titlesec}
\usepackage{url}
\usepackage{textcomp}
\usepackage{xcolor}
\usepackage{listings}
\usepackage{multirow}
\usepackage{longtable}
\usepackage{tabularx}
\usepackage{times}

% Set up page style
\pagestyle{fancy}
\fancyhf{}
\rfoot{\thepage}
\renewcommand{\headrulewidth}{0pt}

% Title formatting
\titleformat{\chapter}[display]{\normalfont\bfseries}{\chaptertitle\thechapter}{0pt}{\Huge}
\titlespacing*{\chapter}{0pt}{50pt}{40pt}

% Document info
\title{Template Document}
\author{Academic Support Agent}
\date{\today}

\begin{document}
\maketitle

\setlength{\parindent}{15pt}


\vspace{1\baselineskip}
\addcontentsline{toc}{section}{Summary}
\section*{Summary}
Summary of report and insights. 


\newpage

\addcontentsline{toc}{section}{Introduction}
\section*{Introduction}
Starts at the top of the next page. The introduction should contain background information on the dataset and community partner, context for the project, and the prompt or desired analyses from the partner.

New paragraphs should be indented but have no space added between lines. A double return should be used at the end of sections to provide visual contrast


\addcontentsline{toc}{section}{Preparation}
\section*{Preparation}

Level 1 headers are Outfit Semibold size 14 font. Body text is Outfit Normal size 12. The document (except for the references) is double spaced.

Use the preparation section to describe the data as your received it and to describe your wrangling, cleaning, and transformation.

\subsection*{Subheader}

Level 2 headers can be added using \begin{verbatim}
\subsection*{name}
\end{verbatim}

\subsubsection*{3rd Level Header}
Level 3 headers can be added using \begin{verbatim} 
\subsubsection*{name}
\end{verbatim}

\addcontentsline{toc}{section}{Analysis}
\section*{Analysis}
Describe the analyses performed. Ensure your descriptions will be understandable to your intended audience. Link to articles/examples in your references if additional explanation would benefit your readers. 

\addcontentsline{toc}{section}{Results}
\section*{Results}
You may choose to combine your results and analysis sections and divide instead by topics if need be.

\addcontentsline{toc}{section}{Recommendations}
\section*{Recommendations}
If this project is unable to address the requests of the partner, use this section to identify future data to be collected which may better suit their inquiries. Acknowledge the limitations of this project and put a positive spin on it.

\addcontentsline{toc}{section}{Conclusion}
\section*{Conclusion}
Answer the community partner (TM)s prompt. Put your results back in the context of the industry
\autocite{Ex95}

\newpage
\begin{singlespace}
\printbibliography

Guidelines and examples of this style are available at: \href{https://owl.purdue.edu/owl/research_and_citation/chicago_manual_17th_edition/cmos_formatting_and_style_guide/general_format.html}{\uline{\textcolor[HTML]{1155CC}{purdue.edu/owl/chicago}}} This section should start on a new page and be single spaced. All references should use hanging indentation 

\end{singlespace}
\end{document}